\NewDocumentCommand{\URL}{m}{\href{#1}{\texttt{#1}}}

% package siunitx -- custom units

\NewDocumentCommand{\errorMessage}{+m}{\colorbox{red}{#1}}
\NewDocumentCommand{\todoMessage}{+m}{\colorbox{orange}{ToDo: #1}}

% Protected input -- \input entire file or typeset a warning message saying that it does not exist
\NewDocumentCommand{\protectedInput}{m}{%
    \IfFileExists{#1}{%
        \input{#1}%
    }{%
        \errorMessage{Missing file \texttt{#1}!}\\
    }%
}

% Try input -- \input entire file if it exists
\NewDocumentCommand{\tryInput}{m}{%
    \IfFileExists{#1}{%
        \input{#1}%
    }{}%
}

\NewDocumentCommand{\cutHere}{}{%
    \noindent%
    \raisebox{-2.8pt}[0pt][0.75\baselineskip]{\small\ding{34}}%
    \unskip{\tiny\dotfill}
}

\NewDocumentCommand{\insertPicture}{O{} m}{%
    \begin{center}
        \IfFileExists{\activeDirectory/#2}{%
            \includegraphics[#1]{\activeDirectory/#2}%
        }{%
            \includegraphics[#1]{example-image}%
        }%
    \end{center}
}

% An answer that is nothing but a picture. Two things went wrong here, and the eleven
% `answer.md` files that opened with an escaped non-breaking space and then a negative
% \vspace -- from -8mm to -13mm, one tuned number per drawing -- were papering over both.
%
% First, a run-in \subsection title needs a paragraph to sit in, and \insertPicture's
% `center` is vertical-mode material, so the title got deferred to the *next* paragraph and
% the problem number came out under its own picture. The non-breaking space was there to
% open a paragraph for the number to land in. \leavevmode does that without putting a
% character on the page.
%
% Second, `center` contributes \topsep above the drawing, which then floated a line and a
% half clear of the number. Here the picture instead stays in the number's own paragraph,
% raised so its top edge is level with the top of that line, and centred between two \hfill
% in whatever width the number leaves -- so it cannot collide with the number however wide
% it grows, and there is no length left to tune when a drawing is rescaled. A picture shorter
% than a line of text has no gap to close and is left on the baseline where it belongs.
%
% \parfillskip has to go to zero for the trailing \hfill to survive the line break, and is
% put back straight after, so an `answer-extra` following the picture still sets ragged.
% Local to a group: only the answer block turns this on, and a picture inside a problem or a
% solution keeps the breathing room \topsep is there for.
\newsavebox{\dgsPictureBox}
\NewDocumentCommand{\tightPictures}{}{%
    \RenewDocumentCommand{\insertPicture}{O{} m}{%
        \sbox\dgsPictureBox{%
            \IfFileExists{\activeDirectory/##2}{%
                \includegraphics[##1]{\activeDirectory/##2}%
            }{%
                \includegraphics[##1]{example-image}%
            }%
        }%
        \leavevmode
        \edef\dgsRestoreParfillskip{\parfillskip=\the\parfillskip\relax}%
        \parfillskip=0pt
        \hfill
        \ifdim\ht\dgsPictureBox>\ht\strutbox
            \raisebox{\dimexpr\ht\strutbox-\ht\dgsPictureBox\relax}{\usebox\dgsPictureBox}%
        \else
            \usebox\dgsPictureBox
        \fi
        \hfill\hbox{}%
        \par
        \dgsRestoreParfillskip
    }%
}

% The join between the answer and an `answer-extra`, `answer-interval` or `answer-also`.
%
% Every part is a separate pandoc file and so ends in a newline, which TeX turns into a
% space; the negative thin space pulls that back, so the comma sits tight against the
% fragment before it. All of that is right only while the answer is running text. An answer
% that is a picture -- or a list, or a display, or anything else block-level -- has already
% ended its paragraph, so the comma opens a new one and prints on a line of its own with
% nothing before it. That is what `21/pv-to-vt-2` was doing: a diagram, then a paragraph
% opening ", Dbajte na to, ...".
%
% \ifvmode asks TeX which of the two actually happened instead of guessing from the source,
% so this is also right for an `answer.md` that is empty -- a problem whose whole answer
% lives in the extra, `28/john-doe` and `20/big-brother` -- and for whatever block-level
% answer is written next.
\NewDocumentCommand{\answerJoin}{}{%
    \ifvmode\noindent\else\!\textit{, }\fi
}

% The trailing block of a booklet -- the four author lists and the colophon under their rule.
% It is one credit, so it reads as one thing or not at all, and `21/sk` was splitting it
% four names from the end: the `Obrázky` heading dangling on the foot of page 48 with its one
% name overleaf, and the colophon stranded at the top of a mostly empty page 49.
%
% \vfill alone cannot prevent that. Glue is a legal breakpoint, so when the block does not
% fit TeX breaks inside it, and the \vfill that was meant to seat it at the foot of the page
% is discarded at the break -- which is why the remainder came out at the *top* of the next
% page rather than the bottom. The fill has to be \vspace*{\fill} for the same reason from
% the other side: plain glue is discarded at the head of a fresh page too, so after the
% \newpage below it would vanish and leave the block at the top again. The star is what makes
% it survive.
%
% So the block is collected into a box first, which cannot be broken at all, and only then is
% there anything to measure: if what is left of the page will not hold it, take a new page
% *before* laying down the \vfill, so the block is seated at the foot of whichever page it
% ends up on. \pagegoal is \maxdimen while a page is still empty, which is not a height and
% must not be compared against one.
%
% The \par\penalty is not decoration. \pagegoal and \pagetotal only describe the page as of
% the last time the page builder ran, and after an answer they can still be describing the
% page before: measured without it, `chem/03` read 744pt used of 692pt on a page holding one
% structural formula, and took a page it did not need. Contributing a penalty exercises the
% builder, and -100 doubles as a mildly attractive breakpoint, so if TeX does have to break
% it breaks here rather than somewhere worse. The idiom is `needspace`'s \Needspace*.
%
% It follows that a block taller than a whole page cannot be placed: it would overrun the
% bottom margin and say so, loudly, as an Overfull \vbox. At four author lists it is nowhere
% near, and splitting it would be the wrong answer anyway.
\newsavebox{\dgsFooterBox}
\NewDocumentEnvironment{footerBlock}{}{%
    \par
    \begin{lrbox}{\dgsFooterBox}%
    \begin{minipage}{\linewidth}%
}{%
    \end{minipage}%
    \end{lrbox}%
    \par\penalty-100
    \ifdim\pagegoal<\maxdimen
        \ifdim\dimexpr\pagegoal-\pagetotal\relax<\dimexpr\ht\dgsFooterBox+\dp\dgsFooterBox\relax
            \newpage
        \fi
    \fi
    \vspace*{\fill}%
    \usebox{\dgsFooterBox}%
}

% \exampleIO[verbatim input][verbatim output]
% Input and output for KSP
%   - verbatim input
%   ₋ verbatim output
% Uses ExplSyntax from LaTeX3
\ExplSyntaxOn
\char_set_catcode_other:n{`\^^M}
\NewDocumentCommand{\exampleIO}{+v +v}{
    \tl_set:Nn \l_tmpa_tl {#1}
    \tl_set:Nn \l_tmpb_tl {#2}
    \tl_replace_all:Nnn \l_tmpa_tl {^^M} {\par}
    \tl_replace_all:Nnn \l_tmpb_tl {^^M} {\par}

    \begin{minipage}[t]{0.48\linewidth}
        \begin{center}vstup\end{center}
        \vspace{-15pt}                                              % one nasty hack here
        \fbox{
            \begin{minipage}[t]{\linewidth}
                \mbox{}\\[-1.5\baselineskip]                        % another nasty hack here
                \ttfamily
                \l_tmpa_tl
            \end{minipage}
        }
    \end{minipage}

    \begin{minipage}[t]{0.0393\linewidth}
        \mbox{}
    \end{minipage}

    \begin{minipage}[t]{0.48\linewidth}
        \begin{center}výstup\end{center}
        \vspace{-15pt}                                              % guess what here
        \fbox{
            \begin{minipage}[t]{0.975\linewidth}
                \mbox{}\\[-1.5\baselineskip]                        % and one nasty hack here
                \ttfamily
                \l_tmpb_tl
            \end{minipage}
        }
    \end{minipage}
    \\[1ex]
}
\char_set_catcode_end_line:n{`\^^M}

\ExplSyntaxOff
